% example-accessibility.tex
% Accessibility demonstration

\DocumentMetadata{}
\documentclass{ltx-talk}
\usepackage{ltxtalk-theme-hawk}

\useltxtalktheme{hawk}

\title{Accessible Presentation}
\author{Accessibility Team}
\institute{University of Inclusive Design}
\date{\today}

\begin{document}

% Die Abschnitts- und Folienstruktur wird von ltx-talk getaggt; das Theme
% platziert und gestaltet sie nur. Explizite Titel statt visueller Umbrüche
% verwenden, damit jede Folie eine echte Überschrift hat.

\section{Introduction to Accessibility}

\begin{frame}
  \frametitle{Introduction to Accessibility}
  \begin{itemize}
    \item What is accessibility?
    \item Why does it matter?
    \item Legal requirements
  \end{itemize}
\end{frame}

\begin{frame}
  \frametitle{Standards and Guidelines}
  \begin{itemize}
    \item WCAG guidelines
    \item PDF/UA standard
    \item Testing tools
  \end{itemize}
\end{frame}

\section{Implementation Strategies}

\begin{frame}
  \frametitle{Implementation Strategies}
  \begin{itemize}
    \item Semantic structure
    \item Alternative texts
    \item Color contrast
  \end{itemize}
\end{frame}

\begin{frame}
  \frametitle{Testing and Feedback}
  \begin{itemize}
    \item Keyboard navigation
    \item Screen reader testing
    \item User feedback
  \end{itemize}
\end{frame}

\section{Best Practices}

\begin{frame}
  \frametitle{Best Practices}
  \framesubtitle{Summary of Key Points}
  \begin{enumerate}
    \item Always provide text alternatives
    \item Ensure sufficient color contrast
    \item Test with actual users
  \end{enumerate}
\end{frame}

\end{document}
